\begin{abstract}
% TODO: 待 Introduction/Method/Experiments 定稿后再写。占位结构（跟模板同款五句式）：
% 1句问题：预测基因扰动下的细胞响应通常靠昂贵的黑盒深度模型，机制不可解释；VCWorld
%   证明"结构化知识检索 + LLM因果推理"这条白盒路线在人类癌细胞系的转录响应预测上可行，
%   但只验证了单一读出层（表达）。
% 1句为什么重要：细胞对扰动的响应本身是多层的（转录、生长/适合度等），这条白盒路线
%   在不同读出层上是否同样有效，是它能否扩展成真正"虚拟细胞"框架的关键未知数——
%   模式生物酵母因为知识图谱（YEASTRACT/STRING/SGA）密度和置信度都高于人类，是检验
%   这个问题最干净的场景。
% 1句方法：把 VCWorld 移植到酵母，构造两个平行任务——T1扰动→表达（Deleteome）、
%   T2A扰动→生长表型（Hillenmeyer 2008），4个本地开源LLM（Qwen2.5-14B/7B、Qwen3-4B、
%   Llama3.1-8B）× 3 seed跑满，统一用macro-F1/MCC（类别不均衡）评测。
% 1-2句关键数字（写作时从 RESULTS.md/task_results.md 摘录当时已定稿的数字，不要提前
%   写死——写这份骨架时 A/B prompt消融、GAT等baseline仍在补充中）：
%   T1a-DE/T1b-DIR 的 macro-F1/MCC 随模型规模单调上升（14B领先），复现VCWorld在人类
%   数据上的"越强越准"；T2A-DIR（敏感/抗）四个模型的MCC全部挤在近随机区间且不随规模
%   变化，是本文的关键负结果。
% 1句意义：说明"知识检索+LLM推理"这条白盒路线对"扰动→调控/表达"类问题有效，但对
%   "扰动→定量剂量依赖的适合度"类问题不够，为虚拟细胞建模指出哪类生物学问题需要
%   补充数值/定量特征而非纯文本证据。
\end{abstract}
